\chapter*{Übersicht von KI-Tools}
\noindent Grundvoraussetzung einer zutreffenden Leistungsbewertung ist, dass der Prüfling die für den Erfolg seiner oder ihrer Prüfung maßgeblichen Leistungen persönlich und ohne fremde Hilfe bzw. sonstiger Hilfsmittel erbringt, soweit diese nicht im Einzelfall zugelassen sind.

\noindent Vor diesem Hintergrund und zum Zwecke der Gewährleistung einer entsprechend eigenständigen Prüfungsleistung sind die in der Anlage aufgeführten, (auch) KI-gestützten, digitalen Tools bei der Bearbeitung der o.g. Prüfung – ggf. mit entsprechender Kenntlichmachung – zugelassen.

\noindent Die Prüflinge tragen auch beim Einsatz von Hilfsmitteln die Verantwortung für die gesamte eingereichte Prüfungsleistung. Dies schließt Erzeugnisse von KI-Tools explizit mit ein. Selbst wenn diese Liste die Verwendung einzelner Tools für zulässig erklärt, sind die Ergebnisse vor einer Übernahme auf ihre sachliche Richtigkeit, Einschlägigkeit und Vereinbarkeit mit den Prüfungsanforderungen hin zu überprüfen.

\noindent\textit{Optional kann für nicht erwähnte Tools eine Freigabe nach Rücksprache zugelassen werden. Hierfür ist dann der nachfolgende Absatz in die Erklärung aufzunehmen. Wird eine Freigabe erteilt, ist die Liste zu aktualisieren, mit einem neuen Datum zu versehen und gegenüber allen betroffenen Studierenden erneut bekannt zu machen, um eine Benachteiligung anderer ausschließen zu können.}

\noindent Angesichts der hohen Geschwindigkeit des technischen Fortschritts auf dem Gebiet der – (auch) KI-gestützten – digitalen Tools kann es dazu kommen, dass die vorstehende Tabelle nicht sämtliche Tools berücksichtigt. Sollten Sie ein Tool bei der Bearbeitung der o.g. Prüfung verwenden wollen, welches in der vorstehenden Tabelle nicht enthalten ist und welches die Eigenständigkeit der o.g. Prüfungsleistung nicht beeinflusst, so halten Sie bzgl. einer etwaigen Verwendung Rücksprache mit der Prüferin / dem Prüfer. Eine Verwendung ohne entsprechende Freigabe durch die Prüferin / den Prüfer ist ausdrücklich nicht gestattet.

\newpage
\begin{longtable}{|p{4.5cm}|p{5.cm}|p{4.5cm}|}
\caption*{Übersicht über die verwendeten KI-Tools für die Abschlussarbeit (Stand: \InfoStateToolsAI)}\\\hline
\textbf{Tool} (ggf. mit technischen Einschränkungen) & \textbf{Zweck der Verwendung} & \textbf{Nachweis- oder Kennzeichnungspflicht} \\\hline\hline
Microsoft Word oder vergleichbare Textverarbeitungsprogramme \newline \textbf{ohne} KI-Plugins/Addins (z.B. Co-Pilot, Writefull) & für alle Zwecke & kein gesonderter Nachweis erforderlich \\
\hline
Google oder vergleichbare Suchmaschinen & für die Recherche & kein gesonderter Nachweis erforderlich \\
\hline
\multirow{6}{=}{ChatGPT und andere LLMs über ChatAI-Zugang der UDE} 
& Materialsuche, Literaturrecherche, -verwaltung, Zusammenfassung, Formulierungsvorschläge, Gedankenanstöße & kein gesonderter Nachweis erforderlich \\
\cline{2-3}
& Konzeption, Strukturierung, Plausibilisierung, redaktionelle Überarbeitung & schriftliche, beschreibende Dokumentation im Anhang \\
\cline{2-3}
& Übernahme von Zitaten, inhaltliche Überarbeitung & Kenntlichmachung nach Zitationsregeln, Fußnote mit Promptangabe im Anhang \\
\cline{2-3}
& Datenaufbereitung, -analyse (auch mit Code), Analyse von Plots & schriftliche, beschreibende Dokumentation im Anhang \\
\cline{2-3}
& Tabellen, Visualisierungen (nur mit Code), Datenverbalisierung & Kenntlichmachung der Tabellen, Grafiken und Verbalisierungen \\
\cline{2-3}
& Code Review, Debugging, Testempfehlungen, Software-Design, Skeletons & schriftliche Dokumentation im Anhang, Quelltext, Kennzeichnung im Code-Kommentarblock \\
\hline
\multirow{2}{=}{Writefull (auch als Word- oder Overleaf-Plugin)} 
& Lektorat, Formulierungsvorschläge, redaktionelle Überarbeitung, LaTeX-Code-Generierung & schriftliche, beschreibende Dokumentation im Anhang \\
\cline{2-3}
& Übernahme von Zitaten, inhaltliche Überarbeitung & Kenntlichmachung nach Zitationsregeln \\
\hline
Github Copilot & Skeletons, Hilfsfunktionen, Boilerplate-Code, Software-Design, Refactoring, Auto-Completion & schriftliche Dokumentation im Anhang, Quelltext \\
\hline
DeepL oder vergleichbare KI-Übersetzungstools & Übersetzung fremder Texte zum Verständnis, Übernahme in den Text & kein gesonderter Nachweis erforderlich \\
\hline
DALL-E oder vergleichbare bildgenerierende KI-Tools & Entwurf von Grafiken ohne Übernahme & Kenntlichmachung der Darstellung oder Grafik \\
\hline
\end{longtable}